% Options for packages loaded elsewhere
\PassOptionsToPackage{unicode}{hyperref}
\PassOptionsToPackage{hyphens}{url}
\PassOptionsToPackage{dvipsnames,svgnames,x11names}{xcolor}
\documentclass[
  10pt,
]{article}
\usepackage{xcolor}
\usepackage[margin=0.68in]{geometry}
\usepackage{amsmath,amssymb}
\setcounter{secnumdepth}{-\maxdimen} % remove section numbering
\usepackage{iftex}
\ifPDFTeX
  \usepackage[T1]{fontenc}
  \usepackage[utf8]{inputenc}
  \usepackage{textcomp} % provide euro and other symbols
\else % if luatex or xetex
  \usepackage{unicode-math} % this also loads fontspec
  \defaultfontfeatures{Scale=MatchLowercase}
  \defaultfontfeatures[\rmfamily]{Ligatures=TeX,Scale=1}
\fi
\usepackage{lmodern}
\ifPDFTeX\else
  % xetex/luatex font selection
\fi
% Use upquote if available, for straight quotes in verbatim environments
\IfFileExists{upquote.sty}{\usepackage{upquote}}{}
\IfFileExists{microtype.sty}{% use microtype if available
  \usepackage[]{microtype}
  \UseMicrotypeSet[protrusion]{basicmath} % disable protrusion for tt fonts
}{}
\usepackage{setspace}
\makeatletter
\@ifundefined{KOMAClassName}{% if non-KOMA class
  \IfFileExists{parskip.sty}{%
    \usepackage{parskip}
  }{% else
    \setlength{\parindent}{0pt}
    \setlength{\parskip}{6pt plus 2pt minus 1pt}}
}{% if KOMA class
  \KOMAoptions{parskip=half}}
\makeatother
\usepackage{longtable,booktabs,array}
\usepackage{caption}
\captionsetup[table]{skip=6pt}
\newcounter{none} % for unnumbered tables
\usepackage{calc} % for calculating minipage widths
% Correct order of tables after \paragraph or \subparagraph
\usepackage{etoolbox}
\makeatletter
\patchcmd\longtable{\par}{\if@noskipsec\mbox{}\fi\par}{}{}
\makeatother
\setlength{\emergencystretch}{3em} % prevent overfull lines
\providecommand{\tightlist}{%
  \setlength{\itemsep}{0pt}\setlength{\parskip}{0pt}}
\usepackage{fvextra}
\RecustomVerbatimEnvironment{verbatim}{Verbatim}{breaklines=true,breakanywhere=true,fontsize=\footnotesize}
\usepackage{etoolbox}
\AtBeginEnvironment{longtable}{\footnotesize}
\setlength{\parskip}{0.35em}
\setlength{\emergencystretch}{3em}
\usepackage{bookmark}
\IfFileExists{xurl.sty}{\usepackage{xurl}}{} % add URL line breaks if available
\urlstyle{same}
% fallback for those not using the hyperref driver hyperxmp:
\makeatletter
\@ifundefined{xmpquote}{\newcommand{\xmpquote}[1]{#1}}{}
\makeatother
\hypersetup{
  pdftitle={HW1 design report: a bounded agent for Bandit alert classification},
  colorlinks=true,
  linkcolor={blue},
  filecolor={Maroon},
  citecolor={Blue},
  urlcolor={blue},
  pdfcreator={LaTeX via pandoc}}

\title{HW1 design report: a bounded agent for Bandit alert
classification}
\author{}
\date{}

\begin{document}
\maketitle

\setstretch{1.04}
Implementation: \texttt{src/hw1\_agent/}. Run instructions:
\url{README.md}. Section numbers follow the rubric. All numbers in this
report are measured from the committed traces in \texttt{runs/}.

\section{Implementation at a glance}\label{implementation-at-a-glance}

{\def\LTcaptype{none} % do not increment counter
\begin{longtable}[]{@{}
  >{\raggedright\arraybackslash}p{(\linewidth - 2\tabcolsep) * \real{0.5000}}
  >{\raggedright\arraybackslash}p{(\linewidth - 2\tabcolsep) * \real{0.5000}}@{}}
\toprule\noalign{}
\begin{minipage}[b]{\linewidth}\raggedright
Piece
\end{minipage} & \begin{minipage}[b]{\linewidth}\raggedright
Choice
\end{minipage} \\
\midrule\noalign{}
\endhead
\bottomrule\noalign{}
\endlastfoot
Language / runtime & Python 3.12, Bandit 1.9.4 (invoked as
\texttt{python\ -m\ bandit} from the same interpreter) \\
Model & OpenAI \texttt{gpt-5.4-mini} via Chat Completions with function
tools, \texttt{reasoning\_effort\ =\ "none"},
\texttt{tool\_choice\ =\ "required"} \\
Agent shape & one ReAct-style loop per alert, written as an explicit
hand-rolled state machine (\texttt{graph.py}), no agent framework \\
Model-facing tools & \texttt{read\_file}, \texttt{search\_code},
\texttt{list\_dir}, plus \texttt{submit\_verdict} to finish an alert \\
Budget & provider-reported
\texttt{prompt\_tokens\ +\ completion\_tokens}, gated before every call
with a byte-based upper bound on the next input \\
Modules & \texttt{cli} \ensuremath{\rightarrow} \texttt{config} \ensuremath{\rightarrow} \texttt{graph} (nodes, budget
gate, failsafes, settlement, crash handling) using \texttt{scanner},
\texttt{workspace}, \texttt{tools}, \texttt{dispatcher},
\texttt{provider}, \texttt{trace}, \texttt{report} \\
\end{longtable}
}

\begin{center}\rule{0.5\linewidth}{0.5pt}\end{center}

\section{A1. Behaviours}\label{a1.-behaviours}

\textbf{For each alert the agent:}

\begin{enumerate}
\def\labelenumi{\arabic{enumi}.}
\tightlist
\item
  Hides every other OWASP \texttt{BenchmarkTest*} case (A7) and builds a
  context message: the Bandit fields (test id, CWE, severity,
  confidence, issue text) plus \ensuremath{\pm}25 numbered lines around the flagged
  line, read from the disposable copy.
\item
  Asks the model to trace the value that reaches the sink back to its
  source. Is it attacker-influenced or constant, trusted or sanitised,
  and on which paths? The model proposes \texttt{read\_file},
  \texttt{search\_code} or \texttt{list\_dir} calls, and the runtime
  executes them.
\item
  Accepts a \texttt{submit\_verdict} only after the runtime validates
  it. The label must be 0 or 1, the rationale must cite at least one
  line, and every evidence range must lie in a file the model may read
  and inside lines the runtime actually showed the model for this alert.
\item
  Settles the alert's report row exactly once (\texttt{classified},
  \texttt{error} or \texttt{budget\_exhausted}) and moves on to the next
  alert.
\end{enumerate}

\textbf{The agent must never:}

\begin{itemize}
\tightlist
\item
  \textbf{Run target code.} Nothing from the target is imported,
  installed or executed. Bandit is a static analyser, invoked as
  \texttt{python\ -m\ bandit} with an argument list (no shell). The
  tools only read bytes.
\item
  \textbf{Write to the target.} All tool access goes to a disposable
  copy. Bandit writes its JSON to a temp directory outside the target.
  There are no write tools.
\item
  \textbf{Invent observations.} Only the dispatcher writes observations
  (tool messages). The model's text is never treated as an observation.
  Evidence must be a subset of lines the runtime displayed, so a verdict
  that cites code the model never saw is rejected.
\item
  \textbf{Read withheld material}: \texttt{.git/},
  \texttt{expectedresults*.csv}, or another OWASP case (A7).
\item
  \textbf{Invent usage.} The budget charges only provider-reported
  counters. A call that fails before usage is returned records no usage
  and no \texttt{model\_call} event.
\end{itemize}

\section{A2. State machine}\label{a2.-state-machine}

\begin{verbatim}
stateDiagram-v2
    [*] --> prepare
    prepare --> discover
    discover --> select: bandit ok
    discover --> stop: bandit failed / provider unusable
    select --> triage: alert pending (runtime)
    select --> finalize: queue empty (runtime)
    triage --> gate
    gate --> classify: checks pass (runtime)
    gate --> settle: turn limit (runtime)
    gate --> stop: max_steps / budget_exhausted (runtime)
    classify --> dispatch: tool calls (model)
    classify --> validate: submit_verdict alone (model)
    classify --> gate: no tool call \ensuremath{\rightarrow} nudge (runtime)
    classify --> settle: provider error (runtime)
    classify --> stop: fatal provider error (runtime)
    dispatch --> gate
    dispatch --> stop: no_progress / max_steps (runtime)
    validate --> settle: accepted (runtime)
    validate --> gate: rejected (runtime)
    validate --> stop: no_progress (runtime)
    settle --> select
    stop --> finalize
    finalize --> [*]
\end{verbatim}

Any uncaught exception, at any node, goes to the crash handler
(\texttt{Agent.\_crash}). It settles every unsettled alert as
\texttt{error}, writes \texttt{report.json}, and writes the single
\texttt{terminal} event (\texttt{fatal\_error}).

\section{A3. Nodes and edges}\label{a3.-nodes-and-edges}

{\def\LTcaptype{none} % do not increment counter
\begin{longtable}[]{@{}
  >{\raggedright\arraybackslash}p{(\linewidth - 8\tabcolsep) * \real{0.2000}}
  >{\raggedright\arraybackslash}p{(\linewidth - 8\tabcolsep) * \real{0.2000}}
  >{\raggedright\arraybackslash}p{(\linewidth - 8\tabcolsep) * \real{0.2000}}
  >{\raggedright\arraybackslash}p{(\linewidth - 8\tabcolsep) * \real{0.2000}}
  >{\raggedright\arraybackslash}p{(\linewidth - 8\tabcolsep) * \real{0.2000}}@{}}
\toprule\noalign{}
\begin{minipage}[b]{\linewidth}\raggedright
Node
\end{minipage} & \begin{minipage}[b]{\linewidth}\raggedright
Does
\end{minipage} & \begin{minipage}[b]{\linewidth}\raggedright
Outgoing edge(s)
\end{minipage} & \begin{minipage}[b]{\linewidth}\raggedright
Kind
\end{minipage} & \begin{minipage}[b]{\linewidth}\raggedright
Decided by
\end{minipage} \\
\midrule\noalign{}
\endhead
\bottomrule\noalign{}
\endlastfoot
\texttt{prepare} & resolve commit, open the workspace, build the
provider & \ensuremath{\rightarrow} \texttt{discover} & fixed & N/A \\
\texttt{discover} & Bandit \texttt{tool\_request}/\texttt{tool\_result},
normalise, sort, assign ids & \ensuremath{\rightarrow} \texttt{select} / \texttt{stop} &
conditional & runtime \\
\texttt{select} & pick the next unsettled alert & \ensuremath{\rightarrow} \texttt{triage} /
\texttt{finalize} & conditional & runtime (\texttt{routing\ triage} /
\texttt{finalize}) \\
\texttt{triage} & per-case withholding, context message,
\texttt{state\_change} & \ensuremath{\rightarrow} \texttt{gate} & fixed & N/A \\
\texttt{gate} & turn limit, then \texttt{max\_steps}, then token budget
& \ensuremath{\rightarrow} \texttt{classify} / \texttt{settle} / \texttt{stop} & conditional &
runtime (\texttt{routing\ turn\_limit} / \texttt{max\_steps} /
\texttt{budget\_exhausted}) \\
\texttt{classify} & one model call with \texttt{tool\_choice="required"}
(forced \texttt{submit\_verdict} on the last turn), \texttt{model\_call}
event with reported usage & \ensuremath{\rightarrow} \texttt{dispatch} / \texttt{validate} &
conditional & \textbf{model} (\texttt{routing\ tools} /
\texttt{submit}) \\
& & \ensuremath{\rightarrow} \texttt{gate} (nudge) / \texttt{settle} (error) / \texttt{stop}
(fatal) & conditional & runtime (\texttt{no\_tool\_call} /
\texttt{error} / \texttt{fatal\_error}) \\
\texttt{dispatch} & run each requested tool through the dispatcher & \ensuremath{\rightarrow}
\texttt{gate} & fixed (default) & N/A \\
& & \ensuremath{\rightarrow} \texttt{stop} & conditional & runtime (\texttt{no\_progress} /
\texttt{max\_steps}) \\
\texttt{validate} & check the verdict & \ensuremath{\rightarrow} \texttt{settle} /
\texttt{gate} / \texttt{stop} & conditional & runtime (\texttt{settle} /
\texttt{reject\_verdict} / \texttt{no\_progress}) \\
\texttt{settle} & write the row, one \texttt{state\_change} with
\texttt{alert\_status} + \texttt{label}, restore withheld case files & \ensuremath{\rightarrow}
\texttt{select} & fixed & N/A \\
\texttt{stop} & settle every remaining alert
(\texttt{budget\_exhausted}, or \texttt{error} for the other stops) & \ensuremath{\rightarrow}
\texttt{finalize} & fixed & N/A \\
\texttt{finalize} & write \texttt{report.json}, then the one
\texttt{terminal} event & end & fixed & N/A \\
\end{longtable}
}

Every conditional edge taken writes exactly one \texttt{routing} event
with \texttt{by}. Fixed and default edges are not traced. The model
chooses only between investigating (\texttt{tools}) and finishing
(\texttt{submit}). Every other transition is the runtime's. The
\texttt{no\_tool\_call} nudge edge is a fallback: since
\texttt{tool\_choice} became \texttt{"required"} (see
\emph{Implementation path}), it fired 0 times across the 200 model calls
of the six runs.

\section{A4. Agent state}\label{a4.-agent-state}

{\def\LTcaptype{none} % do not increment counter
\begin{longtable}[]{@{}
  >{\raggedright\arraybackslash}p{(\linewidth - 8\tabcolsep) * \real{0.2000}}
  >{\raggedright\arraybackslash}p{(\linewidth - 8\tabcolsep) * \real{0.2000}}
  >{\raggedright\arraybackslash}p{(\linewidth - 8\tabcolsep) * \real{0.2000}}
  >{\raggedright\arraybackslash}p{(\linewidth - 8\tabcolsep) * \real{0.2000}}
  >{\raggedright\arraybackslash}p{(\linewidth - 8\tabcolsep) * \real{0.2000}}@{}}
\toprule\noalign{}
\begin{minipage}[b]{\linewidth}\raggedright
Field
\end{minipage} & \begin{minipage}[b]{\linewidth}\raggedright
Reducer
\end{minipage} & \begin{minipage}[b]{\linewidth}\raggedright
Lifetime
\end{minipage} & \begin{minipage}[b]{\linewidth}\raggedright
Writer
\end{minipage} & \begin{minipage}[b]{\linewidth}\raggedright
Model sees?
\end{minipage} \\
\midrule\noalign{}
\endhead
\bottomrule\noalign{}
\endlastfoot
\texttt{alerts} (report rows: status, label, rationale, evidence) & set
once per alert in \texttt{\_settle} (a second write raises) & run &
runtime (\texttt{settle}, \texttt{stop}, crash handler) & no \\
\texttt{queue} (unsettled alerts) & remove on settle & run & runtime &
no \\
\texttt{usage} (\texttt{input}, \texttt{output}) & add, from
provider-reported counters & run & runtime (\texttt{classify}) & no \\
trace \texttt{usage} mirror (incl.~\texttt{cache\_read}) & add, from
written \texttt{model\_call} events only & run & \texttt{Trace.event} &
no \\
\texttt{steps\_used} & add 1 per model call and per tool execution & run
& runtime (\texttt{classify}, \texttt{dispatch}) & no \\
\texttt{no\_progress\_streak} & reset to 0 on progress or settlement,
else +1 per turn & run & runtime & no \\
\texttt{provider\_failures} & +1 per failed call, reset on success & run
& runtime & no \\
\texttt{first\_call\_input} & set once (first reported first-turn input)
& run & runtime & no \\
\texttt{stop} (status, reason) & set once & run & runtime & no \\
\texttt{session.messages} & append & alert & runtime
(system/context/nudges/tool results) + model (its own assistant turns) &
\textbf{yes} \\
\texttt{session.observed} (path \ensuremath{\rightarrow} line numbers shown) & union & alert &
dispatcher / triage & no \\
\texttt{session.seen\_requests} (tool+args fingerprints) & union & alert
& dispatcher & no \\
\texttt{session.turns} & add & alert & runtime & no (except the
final-turn instruction) \\
\texttt{session.last\_input}, \texttt{last\_bytes} & overwrite & alert &
runtime & no \\
\texttt{session.pending\_reply} & overwrite & one turn & runtime
(\texttt{classify}) & it is the model's own reply \\
workspace \texttt{active\_case}, stash & set per alert, restored at
settle & alert & workspace & no \\
\end{longtable}
}

The model never writes state directly. Its only outputs are proposals,
tool calls and verdicts, and the runtime validates each one before
anything changes.

\section{A5. Tools}\label{a5.-tools}

All tools take repo-relative POSIX paths and go through
\texttt{Workspace.resolve} (A7) before touching disk. A refusal comes
back as an observation beginning \texttt{REFUSED:} with a
\texttt{tool\_result\ ok=false\ error=\ldots{}} event. It never raises. The
dispatcher refuses an unknown tool name, arguments that are not a JSON
object, arguments that fail the tool's JSON Schema (wrong type, missing
or extra keys), and an exact repeat of an earlier request for the same
alert. Results are capped at 12,000 characters. On truncation, only the
lines still visible count as observed.

{\def\LTcaptype{none} % do not increment counter
\begin{longtable}[]{@{}
  >{\raggedright\arraybackslash}p{(\linewidth - 8\tabcolsep) * \real{0.2000}}
  >{\raggedright\arraybackslash}p{(\linewidth - 8\tabcolsep) * \real{0.2000}}
  >{\raggedright\arraybackslash}p{(\linewidth - 8\tabcolsep) * \real{0.2000}}
  >{\raggedright\arraybackslash}p{(\linewidth - 8\tabcolsep) * \real{0.2000}}
  >{\raggedright\arraybackslash}p{(\linewidth - 8\tabcolsep) * \real{0.2000}}@{}}
\toprule\noalign{}
\begin{minipage}[b]{\linewidth}\raggedright
Tool
\end{minipage} & \begin{minipage}[b]{\linewidth}\raggedright
Arguments
\end{minipage} & \begin{minipage}[b]{\linewidth}\raggedright
Returns
\end{minipage} & \begin{minipage}[b]{\linewidth}\raggedright
Limits
\end{minipage} & \begin{minipage}[b]{\linewidth}\raggedright
Refusals
\end{minipage} \\
\midrule\noalign{}
\endhead
\bottomrule\noalign{}
\endlastfoot
\texttt{read\_file} & \texttt{path}, \texttt{start\_line\ \ensuremath{\geq}\ 1},
\texttt{end\_line\ \ensuremath{\geq}\ start\_line} & numbered lines +
\texttt{{[}file\ has\ N\ lines;\ showing\ a-b{]}} & \ensuremath{\leq} 200 lines per
call; lines \textgreater{} 400 chars truncated & path refused (absolute,
\texttt{..}, outside the copy, withheld); not a file; binary;
\texttt{start\_line} past EOF;
\texttt{end\_line\ \textless{}\ start\_line} \\
\texttt{search\_code} & \texttt{pattern} (literal, 1--200 chars),
\texttt{path} (default \texttt{.}), \texttt{glob} (file-name pattern,
default \texttt{*.py}), \texttt{case\_sensitive} &
\texttt{path:line:\ text} hits & \ensuremath{\leq} 50 hits; files \textgreater{} 2 MB
skipped; \ensuremath{\leq} 30,000 files scanned; literal match only (no regex, so no
ReDoS) & bad pattern/glob; path refused or missing. Withheld files are
silently skipped. \\
\texttt{list\_dir} & \texttt{path} & sorted entries, dirs end in
\texttt{/} & \ensuremath{\leq} 200 entries & path refused; not a directory. Withheld
entries are not listed. \\
\texttt{submit\_verdict} (finish) & \texttt{label\ \ensuremath{\in}\ \{0,1\}},
\texttt{rationale}, \texttt{evidence{[}1..5{]}} of
\texttt{\{path,start\_line,end\_line\}} & Ends the alert & ranges \ensuremath{\leq} 200
lines; rationale stored up to 3,000 chars & rejected if: called together
with other tools; schema fails; rationale empty or cites no line;
evidence path refused or not a file; any cited line was never shown to
the model \\
\end{longtable}
}

Bandit is the runtime's discovery tool, traced as
\texttt{tool\_request}/\texttt{tool\_result} with
\texttt{tool:\ "bandit"}. The model cannot call it.

\section{A6. Start, stop and
termination}\label{a6.-start-stop-and-termination}

\textbf{Start.} The CLI parses \texttt{-\/-input/-\/-output/-\/-config}.
The config must parse as TOML and validate against
\texttt{config.schema.json}. If it does not, the run never starts and
the trace ends \texttt{config\_invalid}. \texttt{{[}agent{]}} keys
(\texttt{max\_turns\_per\_alert}, \texttt{no\_progress\_limit},
\texttt{request\_timeout\_s}) are optional.

\textbf{Terminal statuses.} Closed and mutually exclusive. Exactly one
is written, as the last event.

{\def\LTcaptype{none} % do not increment counter
\begin{longtable}[]{@{}
  >{\raggedright\arraybackslash}p{(\linewidth - 4\tabcolsep) * \real{0.3333}}
  >{\raggedright\arraybackslash}p{(\linewidth - 4\tabcolsep) * \real{0.3333}}
  >{\raggedright\arraybackslash}p{(\linewidth - 4\tabcolsep) * \real{0.3333}}@{}}
\toprule\noalign{}
\begin{minipage}[b]{\linewidth}\raggedright
Status
\end{minipage} & \begin{minipage}[b]{\linewidth}\raggedright
Trigger
\end{minipage} & \begin{minipage}[b]{\linewidth}\raggedright
Rows still unsettled become
\end{minipage} \\
\midrule\noalign{}
\endhead
\bottomrule\noalign{}
\endlastfoot
\texttt{completed} & every alert settled by the agent loop & N/A \\
\texttt{budget\_exhausted} & budget gate:
\texttt{required\ \textgreater{}\ max\_tokens\ \ensuremath{-}\ used} before a model
call & \texttt{budget\_exhausted}, \texttt{label:\ null} \\
\texttt{max\_steps} &
\texttt{steps\_used\ +\ 1\ \textgreater{}\ max\_steps} before a model
call or a tool execution & \texttt{error} (``run stopped
(max\_steps)'') \\
\texttt{no\_progress} & \texttt{no\_progress\_limit} consecutive model
turns with no new observation: no tool call, all requests refused or
duplicated, or verdict rejected & \texttt{error} \\
\texttt{fatal\_error} & uncaught exception, Bandit failure, unusable
provider (missing key, auth, unknown model, invalid parameter, quota),
or 3 consecutive failed calls & \texttt{error} \\
\texttt{config\_invalid} & config rejected before any scan or model call
& (no rows) \\
\end{longtable}
}

\textbf{Precedence} is fixed by where each check runs. A crash pre-empts
everything. In \texttt{gate}, \texttt{max\_steps} is checked before the
budget. \texttt{no\_progress} is evaluated after a turn's observations.
The first check to fire stops the run, so two statuses can never both
apply.

\textbf{Budget policy.}

\begin{itemize}
\tightlist
\item
  \emph{Charge}: \texttt{used\ =\ \ensuremath{\Sigma}\ (input\ +\ output)} over
  \texttt{model\_call} events, as reported by the provider. For OpenAI,
  \texttt{input\ =\ prompt\_tokens} and
  \texttt{output\ =\ completion\_tokens}, which includes reasoning
  tokens.
  \texttt{cache\_read\ =\ prompt\_tokens\_details.cached\_tokens} is
  recorded and never charged. OpenAI reports no cache-write counter.
  \texttt{report.budget.used\_tokens} is computed from the trace's own
  mirror of the \texttt{model\_call} events, so it equals the trace sum
  by construction.
\item
  \emph{Output allowance}: \texttt{budget.output\_allowance} is sent as
  \texttt{max\_completion\_tokens} on every call.
\item
  \emph{When it is checked}: in \texttt{gate}, before every model call.
  Required = \texttt{projected\_input\ +\ output\_allowance}, and the
  call is made only if \texttt{required\ \ensuremath{\leq}\ max\_tokens\ \ensuremath{-}\ used}.
\item
  \emph{Projection} (gate only, never charged): an upper bound built
  from provider numbers. A BPE token covers at least one byte, so the
  UTF-8 bytes added since the previous call bound the tokens added:
  \texttt{projected\ =\ last\ reported\ input\ +\ bytes\ added\ +\ 16/message}.
  A new alert's first call uses the first reported first-turn input plus
  the new alert message's bytes. Before any report exists, it uses the
  whole request's bytes. Because the projection is an upper bound and
  output is capped by the allowance, \texttt{used\ \ensuremath{\leq}\ max\_tokens}
  always holds.
\item
  \emph{At the cutoff}: a \texttt{routing} event
  \texttt{budget\_exhausted} by \texttt{runtime} records
  \texttt{remaining} and \texttt{required}. The current and all later
  alerts get \texttt{budget\_exhausted} rows (one \texttt{state\_change}
  each). Then \texttt{report.json} is written with
  \texttt{exhausted:\ true}, followed by
  \texttt{terminal\ budget\_exhausted}.
\end{itemize}

\textbf{Per-alert bound.} At most \texttt{max\_turns\_per\_alert} (8)
model calls per alert. Every call is sent with
\texttt{tool\_choice\ =\ "required"}, so each turn is either an
investigation or a verdict. On the last one the runtime forces
\texttt{tool\_choice\ =\ submit\_verdict}. If that verdict is still
rejected, the alert settles \texttt{error} and the run continues. In the
real runs this forced turn was used once
(\texttt{radicale/storage/multifilesystem/lock.py:102}, B602: 7
investigation calls, then the forced verdict).

\section{A7. Workspace and
withholding}\label{a7.-workspace-and-withholding}

{\def\LTcaptype{none} % do not increment counter
\begin{longtable}[]{@{}
  >{\raggedright\arraybackslash}p{(\linewidth - 8\tabcolsep) * \real{0.2000}}
  >{\raggedright\arraybackslash}p{(\linewidth - 8\tabcolsep) * \real{0.2000}}
  >{\raggedright\arraybackslash}p{(\linewidth - 8\tabcolsep) * \real{0.2000}}
  >{\raggedright\arraybackslash}p{(\linewidth - 8\tabcolsep) * \real{0.2000}}
  >{\raggedright\arraybackslash}p{(\linewidth - 8\tabcolsep) * \real{0.2000}}@{}}
\toprule\noalign{}
\begin{minipage}[b]{\linewidth}\raggedright
Withheld
\end{minipage} & \begin{minipage}[b]{\linewidth}\raggedright
Targets
\end{minipage} & \begin{minipage}[b]{\linewidth}\raggedright
Why
\end{minipage} & \begin{minipage}[b]{\linewidth}\raggedright
How
\end{minipage} & \begin{minipage}[b]{\linewidth}\raggedright
Restored
\end{minipage} \\
\midrule\noalign{}
\endhead
\bottomrule\noalign{}
\endlastfoot
\texttt{.git/} & all & history may contain commits after the pin,
including fixes & never copied (\texttt{copytree(ignore=.git)});
\texttt{resolve} refuses any \texttt{.git} component & nothing moved;
the original is untouched \\
\texttt{expectedresults*.csv} & OWASP (rule applies everywhere) & the
answer key & moved from the copy into a sibling stash
(\texttt{repo-withheld/}) for the whole run; \texttt{resolve} refuses
the name & moved back on workspace exit, then the copy is deleted \\
every \texttt{BenchmarkTest*} file except the current case's files
(\texttt{.py} and its
\path{templates/web/*/BenchmarkTest<id>.html})
& OWASP & each case has a safe twin, so the answer is one grep away &
moved to the stash in \texttt{triage} for that alert; \texttt{resolve}
refuses non-active case names & moved back in \texttt{settle}, and on
exit \\
\end{longtable}
}

\textbf{One path check.} \texttt{Workspace.resolve(path)} is the only
route from a tool argument to the file system. It refuses non-strings,
NUL bytes, absolute paths (including drive letters), any \texttt{..}
component, withheld names, and anything whose fully resolved path, with
symlinks followed, lies outside the copy, or whose resolved target is
withheld. \texttt{search\_code} and \texttt{list\_dir} call
\texttt{resolve} on every file or entry they would return, so withheld
files are invisible to them too, not only to direct reads. Discovery
(Bandit) scans the original \texttt{-\/-input} in full, as required, and
never writes to it. \texttt{helpers/} is not withheld.

\begin{center}\rule{0.5\linewidth}{0.5pt}\end{center}

\section{Implementation path and
changes}\label{implementation-path-and-changes}

We built the agent offline first and only then pointed it at the real
model. Each change below was driven by a concrete failure, and the trace
showed what went wrong.

\textbf{1. Offline build.} We implemented and checked the components in
this order: config validation, then the Bandit wrapper (reproducing the
brief's counts of 10 / 21 / 18 at the pinned commits), then the
workspace and path check, the trace and report writers, and finally the
graph. To exercise the graph without spending tokens we drove it with a
scripted model and with a local stand-in for the OpenAI endpoint, and
checked every run against the three schemas and the six trace
invariants. Those development harnesses are not part of the submission.

\textbf{2. Provider: OpenAI, and a budget gate without a tokenizer.} We
use OpenAI because that is the API key we had. Chat Completions has no
free endpoint for counting input tokens, and \texttt{tiktoken} downloads
its encoding files at first use, which adds a network dependency. We
kept the gate dependency-free instead: every BPE token covers at least
one byte, so the UTF-8 bytes added to the conversation since the last
provider-reported input are a hard upper bound on the tokens added (A6).
The bound only gates calls. What is charged is always the provider's own
\texttt{prompt\_tokens\ +\ completion\_tokens}.

\textbf{3. Bandit version check.}
\texttt{python\ -m\ bandit\ -\/-version} prints
\texttt{\_\_main\_\_.py\ 1.9.4}, not \texttt{bandit\ 1.9.4}, so our
first version check failed every scan. We now parse the version from the
first line, whatever the program name, and still refuse anything other
than the configured 1.9.4.

\textbf{4. First real run: \texttt{fatal\_error} with 0 steps.} The
provider could not be built because \texttt{OPENAI\_API\_KEY} was not
set in that shell. Bandit still ran, all 10 Radicale alerts got
\texttt{error} rows, and the trace ended with one
\texttt{terminal\ fatal\_error}. This was an unplanned but real check of
the fatal path.

\textbf{5. \texttt{reasoning\_effort\ =\ "low"} was rejected.} The first
call returned HTTP 400: \emph{``Function tools with reasoning\_effort
are not supported for gpt-5.4-mini in /v1/chat/completions \ldots{} set
reasoning\_effort to `none'\,''}. Because \texttt{provider.py}
classifies a \texttt{BadRequestError} as fatal (it would repeat on every
call), the run stopped after 1 step instead of failing all 10 alerts one
by one. We set \texttt{reasoning\_effort\ =\ "none"} in both configs. We
did not move to the Responses API, because that would have meant
rewriting the provider for a model setting.

\textbf{6. \texttt{no\_progress} on the first full OWASP run \ensuremath{\rightarrow}
\texttt{tool\_choice\ =\ "required"}.} With reasoning off, the model
often answered in prose instead of calling a tool. On
\texttt{BenchmarkTest00012} it needed one nudge. On
\texttt{BenchmarkTest00013} it read the case and
\texttt{helpers/utils.py} and then replied four times in a row with
\texttt{finish\_reason\ =\ stop} and no tool call. The no-progress
detector fired as designed
(\texttt{4\ consecutive\ turns\ without\ progress}), ending the run
after 2 of 18 alerts. That trace was replaced when we re-ran all six
runs, and the event sequence above is taken from it. The fault was in
our request rather than the detector: with
\texttt{tool\_choice\ =\ "auto"} the model is allowed to skip tools, but
in this design every turn must be a tool call or a verdict. We changed
it to \texttt{"required"}. After the change, every one of the 200 model
calls across the six runs ended in a tool call, and the nudge path never
fired. We then re-ran \textbf{all six} runs so that the committed
outputs come from one version of the code.

\textbf{7. Budgets.} \texttt{configs/full.toml} (1,500,000 tokens) was
sized for SGLang, the largest target. It used at most 215,298.
\texttt{configs/low.toml} (50,000) was kept because it cuts every target
partway: Radicale after 3 of 10 alerts, SGLang after 3 of 21, and OWASP
after 15 of 18.

\begin{center}\rule{0.5\linewidth}{0.5pt}\end{center}

\section{C2. Results}\label{c2.-results}

Model: \texttt{gpt-5.4-mini} (\texttt{reasoning\_effort\ =\ "none"}).
Each run's config is the corresponding file in
\texttt{configs/targets/}.

{\def\LTcaptype{none} % do not increment counter
\begin{longtable}[]{@{}
  >{\raggedright\arraybackslash}p{(\linewidth - 12\tabcolsep) * \real{0.1429}}
  >{\raggedright\arraybackslash}p{(\linewidth - 12\tabcolsep) * \real{0.1429}}
  >{\raggedright\arraybackslash}p{(\linewidth - 12\tabcolsep) * \real{0.1429}}
  >{\raggedright\arraybackslash}p{(\linewidth - 12\tabcolsep) * \real{0.1429}}
  >{\raggedright\arraybackslash}p{(\linewidth - 12\tabcolsep) * \real{0.1429}}
  >{\raggedright\arraybackslash}p{(\linewidth - 12\tabcolsep) * \real{0.1429}}
  >{\raggedright\arraybackslash}p{(\linewidth - 12\tabcolsep) * \real{0.1429}}@{}}
\toprule\noalign{}
\begin{minipage}[b]{\linewidth}\raggedright
run
\end{minipage} & \begin{minipage}[b]{\linewidth}\raggedright
alerts
\end{minipage} & \begin{minipage}[b]{\linewidth}\raggedright
label 1
\end{minipage} & \begin{minipage}[b]{\linewidth}\raggedright
label 0
\end{minipage} & \begin{minipage}[b]{\linewidth}\raggedright
null
\end{minipage} & \begin{minipage}[b]{\linewidth}\raggedright
used / max tokens
\end{minipage} & \begin{minipage}[b]{\linewidth}\raggedright
terminal
\end{minipage} \\
\midrule\noalign{}
\endhead
\bottomrule\noalign{}
\endlastfoot
radicale-full & 10 & 6 & 4 & 0 & 161304 / 1500000 & completed \\
radicale-low & 10 & 0 & 3 & 7 & 37278 / 50000 & budget\_exhausted \\
sglang-full & 21 & 13 & 8 & 0 & 215298 / 1500000 & completed \\
sglang-low & 21 & 3 & 0 & 18 & 40941 / 50000 & budget\_exhausted \\
owasp-full & 18 & 10 & 8 & 0 & 47336 / 1500000 & completed \\
owasp-low & 18 & 6 & 9 & 3 & 44075 / 50000 & budget\_exhausted \\
\end{longtable}
}

\textbf{OWASP precision / recall (full run):} TP=6, FP=4, FN=2, TN=6 \ensuremath{\rightarrow}
precision 0.600, recall 0.750.

\textbf{How the labels were used to refine the agent.} On the full OWASP
run the agent labelled 12 of the 18 alerts correctly. We compared each
wrong label with its rationale and trace, and the errors fall into three
groups:

\begin{itemize}
\tightlist
\item
  \emph{A guard was seen but not trusted} (false positives:
  \texttt{00073} B307, \texttt{00075} B102). Both rationales cite the
  check on lines 46-50 that only lets a quoted string literal through,
  yet still treat \texttt{eval}/\texttt{exec} as reachable from the
  cookie. The model took ``the value comes from a cookie'' as enough,
  even though the label definition says a value neutralised on every
  path is a 0.
\item
  \emph{Parser defaults were assumed unsafe} (false positives:
  \texttt{00017} B317 and B318). The model reasoned that a
  cookie-controlled document reaching \texttt{xml.sax} and
  \texttt{minidom} ``without defusing'' means XXE, without checking
  whether those parsers resolve external entities by default.
\item
  \emph{Weak hashes were judged by how the digest is used} (false
  negatives: \texttt{00058}, \texttt{00060}, B324). The model called
  SHA-1 safe because the digest is only written to a file and gates no
  security decision. In \texttt{00060} it also leaned on the
  test/dev-only exemption because the file sits under
  \texttt{testcode/}, which is true of every OWASP case. The benchmark
  marks both as vulnerable.
\end{itemize}

We did \textbf{not} change the prompt in response to these errors. The
safe-use criteria and the ``check every path'' instruction were already
in the system prompt from the first run, and the errors above happened
anyway. They show that the model applied those rules unevenly rather
than that the rules were missing. The labels shaped our conclusions
about where the next changes should go, and we deliberately did not
apply them to these 18 cases, because that would only tune the agent to
the development set while Radicale and SGLang are held out:

\begin{itemize}
\tightlist
\item
  ask the model to confirm that a guard can actually be bypassed before
  labelling a guarded sink 1;
\item
  ask it to check a parser's defaults instead of assuming them;
\item
  apply the test/dev-only exemption only when the code really cannot run
  in production, never because of a directory name;
\item
  state that a weak hash on a reachable path is a finding regardless of
  what the digest is used for.
\end{itemize}

The changes we did make during development were to the runtime and the
request (see \emph{Implementation path}, items 5 and 6), not to the
labelling rules.

\textbf{Other measured behaviour.}

{\def\LTcaptype{none} % do not increment counter
\begin{longtable}[]{@{}
  >{\raggedright\arraybackslash}p{(\linewidth - 10\tabcolsep) * \real{0.1667}}
  >{\raggedright\arraybackslash}p{(\linewidth - 10\tabcolsep) * \real{0.1667}}
  >{\raggedright\arraybackslash}p{(\linewidth - 10\tabcolsep) * \real{0.1667}}
  >{\raggedright\arraybackslash}p{(\linewidth - 10\tabcolsep) * \real{0.1667}}
  >{\raggedright\arraybackslash}p{(\linewidth - 10\tabcolsep) * \real{0.1667}}
  >{\raggedright\arraybackslash}p{(\linewidth - 10\tabcolsep) * \real{0.1667}}@{}}
\toprule\noalign{}
\begin{minipage}[b]{\linewidth}\raggedright
run
\end{minipage} & \begin{minipage}[b]{\linewidth}\raggedright
model calls
\end{minipage} & \begin{minipage}[b]{\linewidth}\raggedright
calls per classified alert (mean / max)
\end{minipage} & \begin{minipage}[b]{\linewidth}\raggedright
tokens per classified alert (mean / max)
\end{minipage} & \begin{minipage}[b]{\linewidth}\raggedright
\texttt{cache\_read} share of input
\end{minipage} & \begin{minipage}[b]{\linewidth}\raggedright
wall time
\end{minipage} \\
\midrule\noalign{}
\endhead
\bottomrule\noalign{}
\endlastfoot
radicale-full & 46 & 4.6 / 8 & 16,130 / 37,867 & 65\% & 45 s \\
sglang-full & 73 & 3.5 / 7 & 10,252 / 29,012 & 45\% & 83 s \\
owasp-full & 28 & 1.6 / 3 & 2,630 / 5,150 & 34\% & 51 s \\
\end{longtable}
}

\begin{itemize}
\tightlist
\item
  \emph{Investigation depth follows the target.} On OWASP, 10 of 18
  alerts were decided from the \ensuremath{\pm}25-line triage context alone, with no
  tool call, because each case is self-contained. Radicale and SGLang
  needed cross-file reads: 5 of their 69 evidence ranges cite a file
  other than the flagged one.
\item
  \emph{The runtime checks fired on real model output.} In
  \texttt{owasp-low}, the first verdict for \texttt{BenchmarkTest00059}
  was rejected because its rationale cited no line. The model
  resubmitted on the next turn and the verdict was accepted. In
  \texttt{sglang-full}, a \texttt{read\_file} on a directory path came
  back as a refusal, and the model continued. The forced final turn was
  used once (A6).
\item
  \emph{The budget gate is safe but conservative.}
  \texttt{used\_tokens\ \ensuremath{\leq}\ max\_tokens} held in every run. At the
  Radicale low cutoff, 12,722 tokens remained but the next call required
  13,624: a projected input of 9,528 plus the 4,096 output allowance.
  The largest output of any call across all six runs was 247 tokens
  (mean 85), so most of the reserved headroom is the output allowance.
  An allowance of about 1,024 would let the same budget cover more
  alerts.
\item
  \emph{Charged input includes cached input.} OpenAI cached 34--65\% of
  the prompt tokens on the full runs (\texttt{cache\_read}). We record
  it and do not charge it separately, but \texttt{prompt\_tokens}
  already includes it, so our charge is the conservative one.
\item
  \emph{Labels are not stable across runs.}
  \texttt{radicale/sharing/files.py:131} (B301) was labelled 1 in
  \texttt{radicale-full} and 0 in \texttt{radicale-low}. Within one run,
  similar pickle-cache sinks also diverged: \texttt{cache.py:109} was
  labelled 1, while \texttt{history.py:53} and \texttt{sync.py:79} were
  labelled 0. We do not fix a sampling seed, and one run per
  configuration cannot separate prompt effects from sampling noise.
\item
  \emph{Low-budget OWASP run:} 15 of 18 alerts were classified before
  the cutoff, giving precision 0.50 and recall 0.50 (TP 3, FP 3, FN 3,
  TN 6, 3 unlabelled).
\end{itemize}

\clearpage
\textbf{Trajectory of one alert} (from
\texttt{runs/radicale-full/trace.jsonl}):

\begin{verbatim}
sequenceDiagram
    autonumber
    participant R as Runtime (graph + dispatcher)
    participant M as Model (classifier)
    participant W as Workspace (disposable copy)
    Note over R,W: radicale-radicale-sharing-files.py-131-B301
    Note over R: route triage (by runtime)
    R->>W: withhold other cases, read context 106-156
    R->>M: alert + context
    Note over M: step 36 in=1621 out=32
    M-->>R: proposes read_file
    Note over R: route tools (by model)
    R->>W: read_file radicale/sharing/files.py lines 1-220
    W-->>R: ok, 11062 chars
    R->>M: observation
    Note over M: step 40 in=4446 out=35
    M-->>R: proposes search_code
    Note over R: route tools (by model)
    R->>W: search_code 'pickle.load' in . (*.py)
    W-->>R: ok, 499 chars
    R->>M: observation
    Note over M: step 44 in=4638 out=32
    M-->>R: proposes read_file
    Note over R: route tools (by model)
    R->>W: read_file radicale/sharing/files.py lines 200-390
    W-->>R: ok, 10659 chars
    R->>M: observation
    Note over M: step 48 in=7222 out=174
    M-->>R: proposes submit_verdict
    Note over R: route submit (by model)
    Note over R: verdict accepted
    Note over R: route settle (by runtime)
    Note over R: settled classified label=1
\end{verbatim}

\clearpage
\section{C3. Limitations and
alternatives}\label{c3.-limitations-and-alternatives}

\textbf{Limitations of this design.}

\begin{itemize}
\tightlist
\item
  One ReAct-style loop per alert, so alerts run strictly one after
  another. Latency is the sum of the per-alert loops: 45--83 s per full
  run here.
\item
  The byte-based budget projection plus a 4,096-token output allowance
  is safe but pessimistic. At the Radicale low cutoff it refused a call
  with 12,722 tokens left, while no call ever produced more than 247
  output tokens.
\item
  \texttt{no\_progress} and \texttt{max\_steps} end the whole run, as
  the brief requires, even when only one alert is stuck. We saw this for
  real: one alert answering in prose ended the first full OWASP run
  after 2 of 18 alerts. \texttt{tool\_choice="required"} and the
  per-alert turn cap with a forced final submit now keep one stuck alert
  from triggering them in ordinary cases.
\item
  Labels are not reproducible from run to run (C2), and with a single
  run per configuration the measured precision and recall carry sampling
  noise.
\item
  \texttt{max\_steps} (600) never fired on a real run. The largest run
  used 126 steps, so this failsafe is shown by its implementation (the
  check in \texttt{gate} and \texttt{dispatch}) rather than by a
  committed real trace.
\item
  The evidence check proves the model \emph{saw} the cited lines, not
  that its reasoning about them is right.
\item
  The context is rebuilt from scratch for each alert. Nothing learned
  about shared helpers carries over.
\end{itemize}

\textbf{Alternatives} (analysed, not built):

{\def\LTcaptype{none} % do not increment counter
\begin{longtable}[]{@{}
  >{\raggedright\arraybackslash}p{(\linewidth - 8\tabcolsep) * \real{0.2000}}
  >{\raggedright\arraybackslash}p{(\linewidth - 8\tabcolsep) * \real{0.2000}}
  >{\raggedright\arraybackslash}p{(\linewidth - 8\tabcolsep) * \real{0.2000}}
  >{\raggedright\arraybackslash}p{(\linewidth - 8\tabcolsep) * \real{0.2000}}
  >{\raggedright\arraybackslash}p{(\linewidth - 8\tabcolsep) * \real{0.2000}}@{}}
\toprule\noalign{}
\begin{minipage}[b]{\linewidth}\raggedright
\end{minipage} & \begin{minipage}[b]{\linewidth}\raggedright
This agent (ReAct loop per alert)
\end{minipage} & \begin{minipage}[b]{\linewidth}\raggedright
Fixed workflow
\end{minipage} & \begin{minipage}[b]{\linewidth}\raggedright
Plan-and-execute
\end{minipage} & \begin{minipage}[b]{\linewidth}\raggedright
Supervisor / worker
\end{minipage} \\
\midrule\noalign{}
\endhead
\bottomrule\noalign{}
\endlastfoot
Shape & model chooses tools until it submits; runtime gates every step &
runtime always does the same steps (context \ensuremath{\pm}N lines \ensuremath{\rightarrow} one classify
call) & planner writes an investigation plan per alert; executor runs
it; optional re-plan & supervisor routes alerts (or sub-questions) to
specialised workers (e.g.~taint tracer, crypto-usage checker), then
aggregates \\
Cost & measured 1.6 (OWASP) to 4.6 (Radicale) calls per alert; grows
with the transcript & lowest: 1 call per alert, fixed size & planner
call + executor calls; re-planning adds more & highest: supervisor calls
+ several workers, each with its own context \\
Latency & sequential turns per alert & lowest, and trivially parallel
across alerts & plan is serial; plan steps can run in parallel & workers
can run in parallel, but the supervisor serialises aggregation \\
Coordination & none: one actor; runtime owns state & none & plan \ensuremath{\leftrightarrow}
execution drift; must decide when to re-plan & hardest: hand-offs,
shared state, conflicting worker verdicts \\
Failure behaviour & local: bad requests become refusals; loops caught by
duplicate refusal, turn cap and no-progress & cannot recover from
missing context (e.g.~taint source in another file), so it fails
silently with a wrong label & a wrong plan wastes the whole budget on
irrelevant reads unless validated; re-planning adds loop risk & one
failing worker can stall aggregation; errors compound across hand-offs;
harder to keep the one-settlement-per-row invariant \\
Fit here & good: alerts need a variable amount of cross-file reading &
good for single-file OWASP cases, weak on Radicale/SGLang & useful for
big repos (SGLang) where reading should be planned & probably overkill
for 10--21 alerts \\
\end{longtable}
}

\end{document}
